\section{模型基本假设}

为在保留实际物理系统核心特征的前提下，提炼出具有普适性、严密性且可计算的数学模型，本文结合赛题给出的技术参数与工程实际，提出以下合理假设：

\begin{enumerate}
    \item \textbf{气象与近地风场准稳态假设}：
    假设在整个应急救援任务执行窗口内（约 2.2 小时），灾区天气处于雨后相对稳定状态，近地风力平缓。忽略微尺度湍流、极端突风及气流对无人机飞行的非线性随机扰动。无人机在各航段的爬升、水平巡航及垂直下降速度均视为名义恒定速度，动力学计算采用定常气动与重力做功方程。
    
    \item \textbf{航路规划与空间净空安全假设}：
    假设无人机在任意两个地理节点之间的水平航迹为空间大圆（球面上最小欧氏距离航段）。为避免山地撞击风险，无人机在航段内保持固定海拔巡航，巡航海拔高度统一取该航段沿途 30 米 DEM 地表最高海拔加上 50 米安全净空裕度（即 $H_{cruise} = \max z_{dem} + 50\text{ m}$）；在调度中心 O01 的作业高度取地面海拔（127.7m），在各受灾服务区的投放作业高度取当地地面海拔加上 30 米（即悬停投放净空）。
    
    \item \textbf{物资离散性与不可分割性假设}：
    假设 80 件应急救援物资均为封装完好、规格固定的标准化离散货箱，在空中运输与地面交接过程中具有不可拆分性。各货箱必须整体装载并完整投递至对应的指定服务区，不考虑拆零装载或跨区调整。
    
    \item \textbf{动力电池性能一致性与无记忆效应假设}：
    假设同机型所配备的共享锂电池组出厂一致性优良，健康状态（State of Health, SOH）保持 100\%，且不存在记忆效应。各电池的放电能耗与当前搭载的实际动态质量严格符合题面物理方程，返航充电过程严格服从官方标定的两阶段等效非线性充电函数。
    
    \item \textbf{无线电信道准确定性传播假设}：
    假设 $f=2400\text{ MHz}$ 频段的电磁波在直视视距（LOS）路径上遵循自由空间路径损耗模型（FSPL）；当连线遭遇 30 米 DEM 地表山脊刺入遮挡时，信号产生恒定的 10 dB 地形阻断附加衰落；忽略密林植被对电磁波产生的微观吸收散射及大气的次级折射。
    
    \item \textbf{任务分区管理独立封闭假设}：
    在第四问的多任务组划分框架下，严格遵守应急指挥预案中的“属地负责、各司其职”原则。假设在整个任务周期内，划归各任务组的运输无人机、共享电池、中继无人机及能源组件不进行跨组调配与支援，各组在物理资源上相互独立封闭。
\end{enumerate}
